\chapter{Conclusion}
\label{chap:conclusion}

This thesis studied how AI agents consolidate a sequence of project prompts
into a master prompt while preserving the user's original requirements. The
main concern is that a generated master prompt may appear clear and complete
while containing details that were never requested, retaining requirements
that were later removed, or silently replacing a settled requirement with an
incompatible one.

The traceability experiment showed that the consolidation instruction has a
clear effect on the result. A general request to summarize or preserve meaning
does not always prevent non-traceable content. The \emph{strict} strategy,
which explicitly requires every requirement to come from the original prompt
history, performs better than \emph{basic} in all tested agent--project
comparisons and better than \emph{loose} in most comparisons. However, 
even this direct rule does not eliminate all non-traceable content.

The \emph{spec-guided} strategy explored a different approach by maintaining
the master prompt throughout the original project session. This reduces the
need to reconstruct the complete history at the end, but the results show a
trade-off. The method improves four of the six comparisons with
\emph{strict}, including all three agents on the simple project, but performs
worse for two agents on the hard project. Maintaining the specification on
every turn introduces repeated decisions about whether a prompt adds, changes,
removes, or conflicts with an existing requirement. These decisions become
more difficult in long and conversational prompt histories.

The requirement-level AlignScore evaluation helps identify candidate
requirements that may need inspection. However, the manual review shows that a
low score does not always represent invented content. It may also result from
different wording or from conversational text that was preserved even though
it is not a final requirement. The metric is therefore useful for comparing
candidates and locating possible problems, but it should be used together with
manual review. It also does not measure requirements that were omitted from
the candidate.

Experiment~2 treats the master prompt as a shared project artifact handed to a
new session. It uses both project histories to demonstrate how later prompts
can conflict with settled requirements without directly requesting that the
earlier decisions be changed.

The proposed governance instruction checks each later prompt against the
settled master prompt. Compatible additions are accepted, clear revisions
update the existing requirement, and unclear conflicts are kept separate for
the user to resolve. The method does not attempt to select the correct value
automatically; its purpose is to prevent the agent from silently turning an
ambiguous instruction into a project decision.

Overall, the findings show that a master prompt should not be treated as a
simple summary generated at the end of a coding session. It is a project
specification derived from a sequence of user decisions. Its requirements
should remain traceable to the original prompts, and later changes should be
checked against what has already been settled.

No single prompting strategy completely removes interpretation errors.
Single-pass consolidation must reconstruct the final state from the full
history, while per-turn consolidation introduces many smaller decisions that
can fail independently. A reliable workflow therefore requires both explicit
traceability rules and human review of unclear changes or conflicts.

The practical implication is that teams using a master prompt as a shared
artifact should maintain it carefully rather than regenerate it as an
uncontrolled summary. Keeping the original prompt history, updating the
current specification as decisions change, and asking the user to resolve
unclear conflicts can better preserve intent across AI-assisted development
sessions.
